% ============================================================
% f2015ahw_q09_views.tex
% 2015 安徽卷（文科）第 9 题 · 三视图 —— 清单 F02 / §2
%
% 来源：.work/phase3-b12/figure-inventory.md §2（源 PDF 内容流逐 path 实测真页 pt）
%       源件 .work/phase3-b12/source/2015安徽文.pdf p1
%
% 坐标换算：画布 [x=1pt, y=1pt]（单位 TeX pt）。局部坐标 = 真页 pt − 原点
%   O = (498.358, 577.549)（清单 §2 几何 bbox 左下角）。真页 pt → 局部 pt 为
%   1:1 平移、不缩放——清单 §0.3 的 0.999622 是 pdftocairo 渲染件的 SVG 观测口径换算，
%   只用于与渲染件对读，不入绘图坐标。渲染层面 1 TeX pt = 0.996264 PDF pt（72/72.27），
%   整图统一 0.37%，属排版口径而非坐标换算误差。
%
% 线宽（清单 §0.4 本卷实测分档）：图形轮廓、尺寸线、尺寸界线 0.398pt
% 虚线：无（本图全部实线）
% 箭头：arrows.meta Triangle 长 3.985pt、宽 2.988pt（清单 §6.3 全卷 31 枚同规格）；
%   箭尖落在清单尖端坐标；线端由 TikZ 按箭长回缩（源内线端再多深入 0.398pt = 一个
%   线宽，被填充三角完全覆盖，不可见）。
% 标签：字号 absolute 10.5pt（汉字 STSong 10.5、数字/√2 CMR/CMSY 10.5）；一律 anchor=base 落在清单实测基线上——
%   清单给的是字框中心：汉字字框 = 全角 em 方框（基线 = 中心 − 0.36em），拉丁/数学字框
%   = 字体 ascent/descent 方框（基线 = 中心 − 0.25em）；两式均与清单另给的基线自洽。
%
% 照源复刻的存疑项：
%    - 俯视图右斜边的两条尺寸界线在源中为逐字节相同的两份路径（#77≡#91、#78≡#92，B-S4），
%      落墨完全重合；本图只画一次。
%    - 四种尺寸体系混用（「1」尺寸的尺寸线就在端线连线上、端线长 8.731pt；「√2」尺寸带
%      8.026pt 垂距偏移、界线长 9.878pt 且穿过尺寸线两侧各伸 4.94pt，B-S9），照源。
%    - √2 的上线（源中为独立 rule #93–#96，长 5.250pt）本图并入 $\sqrt{2}$ 字形自带的横线，
%      不另画 4 条 rule；标签位置按清单字框定基线。
% ============================================================
\begin{tikzpicture}[x=1pt, y=1pt,
  every node/.style={inner sep=0, font=\fontsize{10.5pt}{12.6pt}\selectfont},
  line cap=butt, line join=miter, line width=0.398pt,
  arr/.style={-{Triangle[length=3.985pt, width=2.988pt]}}]

  % ---- 主图形（清单 §2.1）----
  \draw (9.168,135.551) -- (96.476,135.551) -- (52.822,179.205) -- cycle;
  \draw (52.822,135.551) -- (52.822,179.205);
  \draw (140.130,135.551) -- (183.784,135.551) -- (140.130,179.205) -- cycle;
  \draw (9.168,70.070) -- (96.476,70.070) -- (52.822,26.416) -- cycle;
  \draw (52.822,26.416) -- (52.822,70.070);

  % ---- 「1」尺寸（清单 §2.2：端线 + 两段尺寸线 + 两端外向箭头）----
  % 正视图下边 ×2
  \draw (9.168,122.454) -- (9.168,131.185);
  \draw (52.822,122.454) -- (52.822,131.185);
  \draw (96.476,122.454) -- (96.476,131.185);
  \draw (12.754,126.820) -- (22.264,126.820);
  \draw (39.725,126.820) -- (49.235,126.820);
  \draw (56.408,126.820) -- (65.918,126.820);
  \draw (83.379,126.820) -- (92.889,126.820);
  \draw[arr] (12.754,126.820) -- (9.168,126.820);
  \draw[arr] (49.235,126.820) -- (52.822,126.820);
  \draw[arr] (56.408,126.820) -- (52.822,126.820);
  \draw[arr] (92.889,126.820) -- (96.476,126.820);
  % 俯视图上边 ×2
  \draw (9.168,74.436) -- (9.168,83.167);
  \draw (52.822,74.436) -- (52.822,83.167);
  \draw (96.476,74.436) -- (96.476,83.167);
  \draw (12.754,78.801) -- (22.264,78.801);
  \draw (39.725,78.801) -- (49.235,78.801);
  \draw (56.408,78.801) -- (65.918,78.801);
  \draw (83.379,78.801) -- (92.889,78.801);
  \draw[arr] (12.754,78.801) -- (9.168,78.801);
  \draw[arr] (49.235,78.801) -- (52.822,78.801);
  \draw[arr] (56.408,78.801) -- (52.822,78.801);
  \draw[arr] (92.889,78.801) -- (96.476,78.801);
  % 侧视图左边 ×1（竖直）
  \draw (127.033,135.551) -- (135.764,135.551);
  \draw (127.033,179.205) -- (135.764,179.205);
  \draw (131.399,139.137) -- (131.399,148.647);
  \draw (131.399,166.108) -- (131.399,175.618);
  \draw[arr] (131.399,139.137) -- (131.399,135.551);
  \draw[arr] (131.399,175.618) -- (131.399,179.205);
  % 侧视图下边 ×1（水平）
  \draw (140.130,122.454) -- (140.130,131.185);
  \draw (183.783,122.454) -- (183.783,131.185);
  \draw (143.716,126.820) -- (153.226,126.820);
  \draw (170.687,126.820) -- (180.197,126.820);
  \draw[arr] (143.716,126.820) -- (140.130,126.820);
  \draw[arr] (180.197,126.820) -- (183.783,126.820);

  % ---- 「√2」尺寸（清单 §2.3：界线 + 两段共线尺寸线 + 两端外向箭头，斜置）----
  % 正视图左斜边
  \draw (6.985,137.734) -- (0.000,144.719);
  \draw (50.639,181.387) -- (43.654,188.372);
  \draw (6.029,143.762) -- (16.589,154.322);
  \draw (34.050,171.783) -- (44.610,182.343);
  \draw[arr] (6.029,143.762) -- (3.493,141.226);
  \draw[arr] (44.610,182.343) -- (47.147,184.880);
  % 正视图右斜边
  \draw (98.658,137.734) -- (105.643,144.719);
  \draw (55.005,181.387) -- (61.990,188.372);
  \draw (89.055,154.322) -- (99.615,143.762);
  \draw (71.593,171.783) -- (61.033,182.343);
  \draw[arr] (99.615,143.762) -- (102.151,141.226);
  \draw[arr] (61.033,182.343) -- (58.497,184.880);
  % 俯视图左斜边
  \draw (6.985,67.887) -- (0.000,60.902);
  \draw (50.639,24.233) -- (43.654,17.248);
  \draw (16.589,51.298) -- (6.029,61.858);
  \draw (34.050,33.837) -- (44.610,23.277);
  \draw[arr] (6.029,61.858) -- (3.493,64.395);
  \draw[arr] (44.610,23.277) -- (47.147,20.741);
  % 俯视图右斜边
  \draw (98.658,67.887) -- (105.643,60.902);
  \draw (55.005,24.233) -- (61.990,17.248);
  \draw (89.055,51.298) -- (99.615,61.858);
  \draw (71.593,33.837) -- (61.033,23.277);
  \draw[arr] (99.615,61.858) -- (102.151,64.395);
  \draw[arr] (61.033,23.277) -- (58.497,20.741);

  % ---- 标签（清单 §2.4；anchor=base 于实测基线）----
  \node[anchor=base] at (52.822,102.301) {正(主)视图};
  \node[anchor=base] at (161.952,102.271) {侧(左)视图};
  \node[anchor=base] at (52.822,1.471) {俯视图};
  \node[anchor=base] at (30.992,123.436) {$1$};
  \node[anchor=base] at (74.652,123.436) {$1$};
  \node[anchor=base] at (30.992,75.416) {$1$};
  \node[anchor=base] at (74.652,75.416) {$1$};
  \node[anchor=base] at (131.402,153.996) {$1$};
  \node[anchor=base] at (161.952,123.436) {$1$};
  \node[anchor=base] at (20.077,159.856) {$\sqrt{2}$};
  \node[anchor=base] at (79.017,159.856) {$\sqrt{2}$};
  \node[anchor=base] at (20.077,37.626) {$\sqrt{2}$};
  \node[anchor=base] at (79.017,37.626) {$\sqrt{2}$};
\end{tikzpicture}
